\section{Método dos Potenciais}

\textbf{Problema base.} O Dijkstra só funciona com pesos $\ge 0$. Com arestas negativas (e sem ciclo negativo),
queremos trocar os pesos por pesos $\ge 0$ \textbf{sem mudar quais caminhos são os menores}.

\textbf{Potencial.} É um número $h(u)$ escolhido para cada vértice, como uma ``altura''. Cada aresta passa a custar
o peso original mais o quanto se desce de $u$ para $v$:
\[
\boxed{w'(u, v) = w(u, v) + h(u) - h(v)}
\]

\subsection*{Por que os menores caminhos não mudam}

Num caminho $s = v_0 \to v_1 \to \dots \to v_k = t$, os potenciais do meio se cancelam
($+h(v_0) - h(v_1) + h(v_1) - h(v_2) + \dots - h(v_k)$), sobrando só as pontas:
\[
w'(\text{caminho}) = w(\text{caminho}) + h(s) - h(t).
\]
\begin{itemize}
  \item Todo caminho de $s$ a $t$ ganha o \textbf{mesmo} valor $h(s) - h(t)$, então o menor continua sendo o menor.
  \item As distâncias se convertem com $d(s, t) = d'(s, t) - h(s) + h(t)$.
  \item Um ciclo começa e termina no mesmo vértice, então seu peso não muda.
\end{itemize}

\subsection*{Como escolher $h$ para ter todo $w' \ge 0$}

\begin{itemize}
  \item $w'(u, v) \ge 0$ é o mesmo que $h(v) \le h(u) + w(u, v)$ em toda aresta.
  \item Distâncias mínimas a partir de uma origem já cumprem isso: se $\text{dist}(v) > \text{dist}(u) + w(u, v)$,
        chegar em $v$ passando por $u$ seria melhor. Então use $h(v) = $ distância mínima até $v$ a partir de um vértice
        virtual ligado a todos com peso $0$ (na prática: Bellman-Ford começando com $h(v) = 0$ para todo $v$).
  \item As arestas usadas pelos menores caminhos a partir da origem ficam com $w' = 0$.
  \item \textbf{Ciclo negativo impede:} somando $h(v) - h(u) \le w(u, v)$ ao longo de um ciclo, o lado esquerdo dá $0$,
        então precisaria $0 \le w(\text{ciclo})$. Logo existe potencial válido $\iff$ não há ciclo negativo.
\end{itemize}
